\begin{musica}{Sete e Pico}{António Mafra}{Tuna de Engenharia da Universidade do Porto}
\begin{multicols}{2}
\nota{Introdução: C G D7 G}

\ac{G}{No} baile da Dona Ester
Feito a semana pas\ac{C}{sa}da
Foram dar com o chauf\ac{D7}{fer}
A dançar com a cri\ac{G}{a}da.
Dizia ela baixinho:
“Na\ac{G7}{pri}se és besti\ac{C}{al}”.

\nota{Refrão:}
\textbf{[}Eram p’raí sete e \ac{G}{pi}co,
\ac{B7}{Oi}to e coisa, nove e \ac{G}{tal}.\textbf{]} \nota{2x}

Chegou a altura da valsa
E exibiu-se o Osório,
De repente cai-lhe a calça,
Rebentou-lhe o suspensório.
Aflito com a mão nos bolsos
Perante o riso geral.

\nota{Refrão}

A Dona Inês sequiosa
Não resistiu ao whisky,
E p’ra se tornar famosa
Quis ir dançar o \textit{twist}.
Ao dar um jeito partiu-se
A coluna vertebral.

\nota{Refrão}

O Dom José de Vicente,
Que é de São Pedro da Cova,
P’ra mostrar que ainda é valente
Foi dançar a bossa nova.
Escorregou no soalho,
Caiu, foi para o hospital.

\nota{Refrão}

Quando o serviço abundante
No baile se iniciou,
O Dom Grilo num instante
A alface devorou.
Diz-lhe a Locas ao ouvido:
“Pareces um animal!”.

\nota{Refrão}

Faltou a luz e gerou-se
A confusão natural,
E a Locas encontrou-se
Nos braços do Amaral.
Logo esta grita aflita:
“Acendam o castiçal!”.

\nota{Refrão}
\end{multicols}

\posicoes{C: x32010, D7: xx0212, G: 320033, G7: 320001, B7: 021202}
\end{musica}
